\subsection{Thematische Einbettung}

\coc\ bettet seine Mechaniken und sein Gameplay in einen immersiven thematischen Kontext und verfügt über eine Story, welche der Spieler im Laufe des Spiels erlebt. 
Das Spiel ist in einer bereits vor einiger Zeit selbst entwickelten Classic-Fantasy-Welt angesiedelt, deren Geschichte mehrere Zeitalter umfasst. 
Die zentralen Spielelemente besitzen dabei nicht nur eine mechanische Funktion, sondern zugleich einen nachvollziehbaren Platz innerhalb dieser Welt. 
Die Level, die der Spieler bewältigt, stellen konkrete Begegnungen, Aufgaben oder Situationen dar, auf die er während seiner Reise trifft. 
Diese können sowohl an bestimmten Orten stattfinden als auch auf dem Weg von einem Ort zum nächsten. 
Die bereisten Orte besitzen einen konkreten Platz und eine Funktion innerhalb der Spielwelt. 
Sie sind auf einer Weltkarte dargestellt und haben thematisch dazu passende Hintergründe während des Levels.
\todo{
    Hier könnte sich evtl. ein Bild anbieten.
}
Auch die eingesetzten Fähigkeiten sind innerhalb der Spielwelt verankert. Sie stellen Zauber dar, die der Spieler beispielsweise von bestimmten Lehrmeistern erlernt und im Verlauf seiner Reise weiterentwickelt. 
Ebenso sind die Charaktere, denen der Spieler begegnet und mit denen er Dialoge führt, in der Geschichte und Struktur dieser Welt plausibel verortet. 
Auf diese Weise soll nicht lediglich eine thematische Kulisse für die Spielmechaniken entstehen. 
Vielmehr soll der Spieler nach und nach ein tatsächliches Verständnis dafür entwickeln, wo er sich befindet, welchen Weg er zurücklegt, wem er begegnet, welche Bedeutung diese Personen und Orte besitzen und weshalb die einzelnen Spielelemente innerhalb dieser Welt existieren.

\todo{Ich bin von `Zeit-Ressource' als Formulierung nach wie vor nicht überzeugt. Vlt. finde ich was Besseres.}

Eine besondere Rolle spielt dabei die Zeit-Ressource, welche der Spieler insbesondere beim \emph{Rasten} innerhalb der Level ausgibt. 
Sie ist sowohl eine zentrale mechanisch-strategische Ressource als auch ein konkreter Bestandteil der Spielwelt. 
Die während eines Levels ausgegebene Zeit stellt dar, wie lange eine Reise zwischen zwei Orten dauert, wie viel Zeit eine Begegnung beansprucht oder wie lange sich die Spielfigur an einem bestimmten Ort aufhält. 
Die Bewegung durch die Spielwelt, das Bewältigen von Levels sowie die Entscheidung, welche Ziele innerhalb eines Runs verfolgt werden, sind damit unmittelbar mit dem verfügbaren Zeitkontingent verbunden. 
Zeit ist folglich nicht lediglich ein abstrakter Zähler oder eine künstliche Begrenzung des Spielsystems, sondern entspricht thematisch der Zeit, die auch innerhalb der Reise der Spielfigur tatsächlich vergeht.

Die Rahmenhandlung erklärt darüber hinaus zentrale Bestandteile der Roguelite-Struktur selbst. 
Insbesondere wird narrativ begründet, weshalb ein Scheitern nicht das endgültige Ende der Reise darstellt, sondern zu einem erneuten Beginn führt.
Die Spielfigur befindet sich innerhalb eines durch die Handlung etablierten Zeitkreislaufs und verfügt in jedem Run nur über ein begrenztes Zeitkontingent, bevor sie an einen Ausgangspunkt -- einen bestimmten Zeitpunkt innerhalb der Handlung -- zurückversetzt wird. 
Im Verlauf des Spiels kann dieses Zeitkontingent erweitert werden. 
Wiederholte Runs, Scheitern, erneute Versuche und der Umgang mit begrenzter Zeit existieren dadurch nicht nur als abstraktes mechanisches Konstrukt, sondern besitzen eine thematisch nachvollziehbare Bedeutung.

Die inhaltliche und thematische Einbettung wird durch die visuelle Präsentation des Spiels unterstützt. Unterschiedliche Orte, Figuren, Levels, Karten und magische Elemente sollen nicht lediglich als funktionale Oberfläche der Spielmechaniken dienen, sondern die Fantasy-welt gestalterisch vermitteln und ihre jeweiligen Charakteristika sichtbar machen. Die visuelle Gestaltung soll dabei helfen, die Immersion zu fördern und die Welt für den Spieler zum Leben zu erwecken. 
Handlung, Weltgestaltung und visuelle Darstellung ergänzen einander somit und können gemeinsam ein verständlicheres und immersiveres Spielerlebnis schaffen, ohne die Übersichtlichkeit und Lesbarkeit des strategischen Kartenspiels zu beeinträchtigen, sondern es zugänglicher und intuitiver zu gestalten.
Die thematische und narrative Ebene ist dabei bewusst kein Ersatz für die mechanische Motivation und nicht der Hauptfokus des Spiels. 
Die grundlegende Spielerfahrung soll bereits durch die strategischen Systeme, das Lösen schwieriger Spielsituationen und die Entwicklung unterschiedlicher Strategien tragen. 
Spielwelt, Geschichte und visuelle Präsentation ergänzen diese Erfahrung jedoch um Bedeutung, Perspektive und Dringlichkeit: Spieler bewältigen eine Herausforderung zum einen, weil deren Lösung spielerisch befriedigend ist, zugleich aber auch, weil sie verstehen, was diese Herausforderung innerhalb der Welt bedeutet und welchem übergeordneten Ziel sie dadurch näherkommen.

\subsubsection{Lore und Handlung}
Die Welt von \coc\ verfügt über eine umfangreichere, mehrere Zeitalter umfassende Hintergrundgeschichte. 
Die folgende Darstellung beschränkt sich bewusst auf die für das Spiel unmittelbar relevanten Zusammenhänge. Die Welt ist darüber hinaus als langfristiger erzählerischer Rahmen angelegt, in dem perspektivisch auch weitere Spiele angesiedelt werden können.
In einem früheren Zeitalter waren die verschiedenen Völker der Welt von Konflikten geprägt. Magie war zu dieser Zeit kaum zugänglich und konnte im Wesentlichen nur von einem Elfenvolk genutzt werden, dessen gewaltiger Baum des Lebens als Quelle magischer Energie diente. Durch die Zusammenarbeit von Wissenschaftlern und Zwergen wurde schließlich entdeckt, dass bestimmte Edelsteine diese Energie speichern können. Auf dieser Grundlage entstand eine Art magischer Kompass, mit dessen Hilfe weitere, bis dahin nicht nutzbare Magievorkommen auf der Welt lokalisiert wurden.
Um diese Quellen zugänglich zu machen, schlossen sich zuvor verfeindete Völker zusammen und pflanzten an den entsprechenden Orten Setzlinge des Baums des Lebens. Aus ihnen entstanden neue magische Bäume, die nach und nach ein weltumspannendes Netz bildeten und Magie für große Teile der Bevölkerung verfügbar machten.
Im folgenden Zeitalter wurde Magie systematisch erforscht, gelehrt und zunehmend in den Alltag der verschiedenen Völker integriert. Es entstanden bedeutende Zentren magischer Forschung und Ausbildung. An ihrer Spitze stehen die Great Sages, die mächtigsten Magier ihrer jeweiligen Disziplin und zugleich zentrale Autoritäten dieser neuen magischen Ordnung.
Zu Beginn von Cycle of Cinder tritt die Spielfigur bei einem dieser Great Sages in die Lehre: einem Magier, der für seine Beherrschung der Zeit bekannt ist. Bereits während der Ausbildung wird Zeitmanipulation als Bestandteil der Welt und der Magie etabliert. So können beispielsweise Trainingssituationen wiederholt werden, indem der Meister die Zeit um einen kurzen Zeitraum zurücksetzt.
Kurz darauf werden Meister und Schüler von einer fremden Entität angegriffen. Dabei erfährt der Spieler, dass zeitgleich auch die übrigen Great Sages ausgeschaltet werden sollen. Ziel dieses koordinierten Angriffs ist es, den wichtigsten Widerstand gegen eine bevorstehende Invasion aus einer parallel existierenden Welt zu beseitigen – jener Welt, aus der die Magie ursprünglich stammt und in der dämonische Wesen leben.
Um seinem Schüler eine Möglichkeit zu geben, diese Ereignisse noch zu verhindern, setzt der Meister einen uralten Zauber ein. Dabei opfert er sich selbst, zerfällt zu Asche und bindet den Spieler an einen Zeitkreislauf: den namensgebenden Cycle of Cinder.
Von diesem Moment an hat der Spieler innerhalb jedes Runs nur ein begrenztes Zeitkontingent, um die Great Sages zu finden und zu retten. Läuft diese Zeit ab oder scheitert der Spieler, wird er an einen früheren Zeitpunkt zurückversetzt und beginnt den Versuch erneut. Im Verlauf des Spiels kann der Spieler seine Möglichkeiten innerhalb dieses Kreislaufs erweitern und unter anderem mehr Zeit für zukünftige Runs gewinnen.
Damit erklärt die Rahmenhandlung unmittelbar zentrale Roguelite-Mechaniken des Spiels: das wiederholte Beginnen eines Runs, das Fortbestehen von Wissen und Fortschritt über mehrere Versuche hinweg sowie Zeit als zentrale Ressource sind nicht nur abstrakte Spielregeln, sondern ergeben sich direkt aus der Handlung und der Logik der Spielwelt.
